\scanpage{53}
\mspage{13}
% archive foliation: 52
\deleted{\unsure{12}}
\uline{Caractérisation galoisienne des filtrations},
\margin{\textcircled{\unsure{B|13}}}

On peut définir le sous-groupe (en fait, anneau) $M_\ell^+(X)$ de $K$ ($K$ des \deleted{\unsure{$\mathbf{Q}_\ell$-faisc.}} $\mathbf{Q}_\ell$-faisceaux constants tordus sur $X$) \uline{provenant} de motifs, on a donc un hom surjectif
\diagram{\begin{tikzcd}
M^+(X) \arrow[r] & M_\ell^+(X)
\end{tikzcd}}
\inserted{Si $X$ est de type fini sur $\mathbf{Z}$,} Cet hom est en fait \uline{bijectif}, car il résulte des conjectures de Tate, plus précisément, que si $E$ et $F$ sont des motifs constants tordus, \uline{simples} \inserted{\uline{non isom.}}, sur $X$, alors $T_\ell(E)$ et $T_\ell(F)$ \deleted{\unsure{ne}} n'ont pas de facteurs simples isomorphes communs.

Par suite, \deleted{\unsure{donc}} la \deleted{\unsure{filtr.}} \uline{graduation} de $M^+(X)$ (i.e. une partition de $\Sigma^+(X)$) et la filtration de $M_\ell(X)$ (i.e. une filtration de $\Sigma_\ell^+(X)$) peuvent se décrire en termes galoisiens. On trouve \inserted{pour $E\in\mathrm{Ob}\,M(X)$} :
\[
\deleted{\mathrm{cl}(E)\text{ de degré }i\ \in}\inserted{\mathrm{cl}(E)\in}M_i(X)
\Longleftrightarrow
\begin{gathered}
\text{pour }\forall\text{ pt fermé }x\text{ de }X,\\
\text{les valeurs propres de Frobenius}\\
\text{dans }T_\ell(E)(\unsure{\bar{x}|x})\text{ sont}\\
\text{de val. abs. }N(x)^{i/2}
\end{gathered}
\]
\margin{``Conjectures de Weil \unsure{et} Riemann''}
